\section{Summary of Findings}
\label{sec:eval-summary}

The empirical findings can be summarised as follows. On the user-experience
instruments (RQ1), Form receives the highest SUS scores per item, the lowest
R-TLX subscale scores, and the largest share of first-place preference votes
($53.6\,\%$). The omnibus tests on SUS and the R-TLX composite do not reach
the conventional significance threshold; the significant sub-effects on the
matched-pairs subset are confined to Form versus the LLM-driven modes on
Temporal Demand, Frustration, and re-use intent. Ask and Chat are
descriptively and inferentially indistinguishable from each other on all
aggregate user-experience measures.

On the per-item Perceived AI Support responses, transparency (item~$3_{ai}$)
is rated highest in Ask, perceived speed-up (item~$2_{ai}$) is rated highest
in Form, and intent to re-use the mode (item~$6_{ai}$) is highest in Form.
Free-text comments characterise Form as fast and ergonomic, Ask as
transparent but cumbersome, and Chat as comfortable but opaque.

On the document-quality side (RQ2), the reference-based NLP metrics place
Form first on lexical recall (BERTScore Recall $0.923$ versus Ask $0.908$
and Chat $0.913$) and on the n-gram-based metrics, Chat first on semantic
similarity (SBERT $0.990$ versus $0.986$ for both other modes) and BLEU,
and Ask consistently last on every metric. The chapter-level breakdown
attributes Chat's BLEU and METEOR advantage primarily to the Retention
Periods and the Technical and Organisational Measures chapters. The
LLM-as-Judge inverts the BERTScore Recall ordering: Form leads on
Completeness ($4.84$) but lags on Scenario Faithfulness ($1.84$) and
Hallucination ($1.65$, inverse-coded), while Chat leads on the latter two
dimensions but trails on Completeness; Overall scores are clustered tightly
between $2.41$ and $2.62$ on the $1$--$5$ scale.
